% Declarações exigidas pelo template da SBC Reviews (versão em português; a versão em inglês está em
% declaracoes-en.tex). Redigidas por Claude Code (Claude Fable 5.2) em 24/11/2026 a partir do que o
% repositório documenta (D143). São declarações dos autores: dependem da confirmação dos três.
\begin{declarations}

\begin{contributions}
% Rascunho segundo a taxonomia CRediT, a confirmar pelos autores.
CCF definiu as questões de pesquisa, os critérios e o protocolo, conferiu a busca, triou as amostras de calibração,
conferiu as decisões de leitura do texto completo, a amostra de extrações e as afirmações empíricas do texto, e
responde pelo conteúdo (conceituação, metodologia, investigação, validação e curadoria de dados). EMLF orientou o
trabalho e PSCG o coorientou (supervisão). A participação das ferramentas de IA está descrita na declaração de uso
de inteligência artificial.
\end{contributions}

\begin{interests}
\pendencia{declaração de conflitos de interesse, a ser feita pelos três autores.}
\end{interests}

\begin{materials}
O protocolo e suas alterações, as consultas exatas de cada base, os prompts, os dados, as fichas de extração, os
scripts e os registros de decisão e de uso de IA criados nesta revisão estão no repositório público do projeto
(\url{https://github.com/caiofreitas1/tcc2-prompt-engineering}). A tabela de apoio e as comparações extraídas podem
ser consultadas no sistema Prompt em Código (\url{https://caiofreitas1.github.io/tcc2-prompt-engineering/}). Os PDFs
e os textos extraídos dos artigos não fazem parte do repositório.
\end{materials}

\begin{aitools}
Este trabalho usou ferramentas de IA generativa na operação da revisão e na redação, como detalha a
Seção~\ref{sec:ia}. Claude Code (Anthropic), com os modelos Claude Opus 5.5, Claude Sonnet 5.5 e Claude Fable 5.2,
executou a busca, a deduplicação, a triagem e a leitura do texto completo, a extração e a avaliação de qualidade, a
síntese, as contagens, a composição das tabelas e o rascunho do texto. O Codex CLI (OpenAI) participou da extração
e auditou materiais do projeto. A conferência dos metadados das referências, o sistema Prompt em Código e a
adaptação a este formato também foram feitos com essas ferramentas. As tabelas e as contagens são geradas por
scripts a partir dos dados do repositório. As etapas conferidas pelo autor e o alcance de cada conferência estão na
Seção~\ref{sec:conferencia}. Os autores respondem por todo o conteúdo.
\end{aitools}

\begin{furtherinformation}
Esta é uma revisão de estudos publicados. Não envolveu participantes humanos nem dados pessoais e não foi submetida
a comitê de ética.
\pendencia{antes da submissão: ORCID e e-mail institucional de Caio César de Freitas; e-mail de Pedro Saint Clair
Garcia; confirmação dos ORCID e do e-mail obtidos em fontes públicas; escolha do autor correspondente; aprovação do
título em inglês; confirmação das contribuições e da leitura da versão final pelos três autores; financiamento, se
houver; e depósito dos artefatos com identificador persistente e licença, recomendado pela revista.}
\end{furtherinformation}

\end{declarations}
